\documentclass[10pt,xcolor={dvipsnames}]{beamer}

% Tema profesional estándar - NO requiere instalación adicional
\usetheme{Copenhagen}
\usecolortheme{seahorse}

% Personalización de colores
\setbeamercolor{structure}{fg=NavyBlue}
\setbeamercolor{palette primary}{bg=NavyBlue,fg=white}
\setbeamercolor{palette secondary}{bg=NavyBlue!80,fg=white}
\setbeamercolor{palette tertiary}{bg=NavyBlue!60,fg=white}
\setbeamercolor{frametitle}{bg=NavyBlue,fg=white}
\setbeamercolor{block title}{bg=NavyBlue!80,fg=white}
\setbeamercolor{block body}{bg=NavyBlue!10,fg=black}

% Configuración adicional
\setbeamertemplate{navigation symbols}{}
\setbeamertemplate{footline}[frame number]

%-------------------------------------------------------
% INCLUDE PACKAGES
%-------------------------------------------------------

\usepackage[utf8]{inputenc}
\usepackage[spanish,es-tabla]{babel}
\usepackage[T1]{fontenc}
\usepackage{graphicx}
\graphicspath{{img/}}
\usepackage{amsmath}
\usepackage{amssymb}

%-------------------------------------------------------
% DEFFINING AND REDEFINING COMMANDS
%-------------------------------------------------------

% colored hyperlinks
\usepackage{hyperref}
\hypersetup{
    colorlinks=true,
    linkcolor=NavyBlue,
    urlcolor=NavyBlue
}

%-------------------------------------------------------
% INFORMATION IN THE TITLE PAGE
%-------------------------------------------------------

\title[]
{
      \textbf{Diseño e Implementación de una Plataforma de Estabilización para una Bicicleta Utilizando un Volante de Inercia}
}

\subtitle[Estabilización Giroscópica]
{
      \textbf{Sistema de Control CMG}
}

\author[Hernández \& Reyes]
{      Jeferson Hernán Hernández Moreno \\
      Fabián Andrés Reyes Romero\\[0.5em]
      {\small Director: Dr. Alexander Jiménez Triana}
}

\institute[]
{%
      Universidad Distrital Francisco José de Caldas\\
      Facultad Tecnológica\\
      Ingeniería en Control y Automatización
}

\date{Enero 2026}

%-------------------------------------------------------
% THE BODY OF THE PRESENTATION
%-------------------------------------------------------

\begin{document}

%-------------------------------------------------------
% THE TITLEPAGE
%-------------------------------------------------------

\begin{frame}[plain]
  \titlepage
\end{frame}


\begin{frame}{Contenido}{}
\tableofcontents
\end{frame}

%-------------------------------------------------------
\section{Introducción}
%-------------------------------------------------------
\begin{frame}{Introducción}{Contexto y Propósito}
%-------------------------------------------------------

\begin{block}{Problemática}
\begin{itemize}
    \item Movilidad sostenible en Bogotá: la bicicleta como alternativa
    \item Inestabilidad crítica a bajas velocidades o en reposo
    \item Riesgo de caídas y accidentes en ciclistas
\end{itemize}
\end{block}

\begin{block}{Propósito del Proyecto}
Desarrollar una plataforma experimental que aproveche el \textbf{efecto giroscópico} para estabilizar activamente una bicicleta en condiciones de reposo mediante un \textbf{volante de inercia} (CMG).
\end{block}

\begin{center}
\fbox{\parbox{0.65\textwidth}{\centering \textit{[Imagen: Sistema CMG en bicicleta]}}}
\end{center}
\end{frame}

%-------------------------------------------------------
\section{Objetivos}
%-------------------------------------------------------
\begin{frame}{Objetivos del Proyecto}{}
%-------------------------------------------------------

\begin{block}{Objetivo General}
Diseñar y construir una plataforma experimental con volante de inercia, sensores y actuadores que permita implementar y validar estrategias de control para la estabilización de una bicicleta.
\end{block}

\begin{block}{Objetivos Específicos}
\begin{enumerate}
    \item Diseñar y construir la estructura mecánica (bicicleta + CMG)
    \item Integrar sistema de sensado (IMU), actuación y procesamiento
    \item Implementar control PID en lazo cerrado y validar experimentalmente
\end{enumerate}
\end{block}

\begin{center}
\fbox{\parbox{0.6\textwidth}{\centering \textit{[Diagrama: Componentes del sistema]}}}
\end{center}
\end{frame}

%-------------------------------------------------------
\section{Marco Teórico}
%-------------------------------------------------------
\begin{frame}{Antecedentes}{Estado del Arte}
%-------------------------------------------------------

\begin{block}{Desarrollos Nacionales}
\begin{itemize}
    \item \textbf{U. Nacional (2014):} Control por dirección, requiere velocidad mínima
    \item \textbf{U. San Buenaventura (2018):} CMG simulado, sin validación física
\end{itemize}
\end{block}

\begin{block}{Desarrollos Internacionales}
\begin{itemize}
    \item \textbf{Ohio State U. (Kalouche, 2014):} CMG experimental, 30s de estabilización
    \item \textbf{Yetkin et al. (2014):} Comparación PID vs SMC
\end{itemize}
\end{block}

\begin{alertblock}{Contribución de este Proyecto}
\textbf{Primera implementación física nacional} de un CMG para estabilización de bicicletas en reposo con validación experimental completa.
\end{alertblock}
\end{frame}

%-------------------------------------------------------
\begin{frame}{Principio Físico}{Efecto Giroscópico}
%-------------------------------------------------------

\begin{columns}
\column{0.5\textwidth}
\begin{block}{Momento Angular}
\begin{equation*}
L = I\omega
\end{equation*}
\end{block}

\begin{block}{Precesión}
\begin{equation*}
\Omega_p = \frac{\tau}{L} = \frac{mgr}{I\omega}
\end{equation*}
\end{block}

\textbf{Principio:} Al inclinar el volante, se genera un torque perpendicular que estabiliza la bicicleta.

\column{0.5\textwidth}
\begin{center}
\fbox{\parbox{0.9\textwidth}{\centering \textit{[Diagrama: Efecto giroscópico]}}}

\vspace{0.3cm}

\fbox{\parbox{0.9\textwidth}{\centering \textit{[Imagen: Volante girando]}}}
\end{center}
\end{columns}
\end{frame}

%-------------------------------------------------------
\begin{frame}{Componentes}{Sensores y Microcontrolador}
%-------------------------------------------------------

\begin{columns}
\column{0.5\textwidth}
\begin{block}{MPU9250 - IMU}
\begin{itemize}
    \item Acelerómetro 3 ejes
    \item Giroscopio 3 ejes
    \item Magnetómetro 3 ejes
    \item Comunicación I²C
\end{itemize}
\end{block}

\begin{center}
\fbox{\parbox{0.8\textwidth}{\centering \textit{[Imagen: MPU9250]}}}
\end{center}

\column{0.5\textwidth}
\begin{block}{ESP32}
\begin{itemize}
    \item Dual-core 32 bits
    \item Hasta 240 MHz
    \item Procesamiento en tiempo real
    \item Control PWM para actuadores
\end{itemize}
\end{block}

\begin{center}
\fbox{\parbox{0.8\textwidth}{\centering \textit{[Imagen: ESP32]}}}
\end{center}
\end{columns}
\end{frame}

%-------------------------------------------------------
\begin{frame}{Controlador PID}{Estrategia de Control}
%-------------------------------------------------------

\begin{block}{Ecuación del Controlador}
\begin{equation*}
u(t) = K_p e(t) + K_i \int_{0}^{t} e(\tau)d\tau + K_d \frac{de(t)}{dt}
\end{equation*}
\end{block}

\begin{columns}
\column{0.5\textwidth}
\begin{itemize}
    \item $K_p$: Acción proporcional (respuesta inmediata)
    \item $K_i$: Acción integral (elimina error residual)
    \item $K_d$: Acción derivativa (amortiguamiento)
\end{itemize}

\column{0.5\textwidth}
\begin{center}
\fbox{\parbox{0.9\textwidth}{\centering \textit{[Diagrama de bloques PID]}}}
\end{center}
\end{columns}

\vspace{0.3cm}
\begin{block}{Filtro Complementario}
\begin{equation*}
\theta = 0.99 \cdot (\theta_{prev} + \omega \cdot \Delta t) + 0.01 \cdot \theta_{accel}
\end{equation*}
Combina precisión del giroscopio con estabilidad del acelerómetro.
\end{block}
\end{frame}

%-------------------------------------------------------
\section{Modelo Matemático}
%-------------------------------------------------------
\begin{frame}{Modelo Dinámico}{Sistema Bicicleta-CMG}
%-------------------------------------------------------

\begin{block}{Componentes del Sistema}
\begin{itemize}
    \item \textbf{Bicicleta ($b$):} Estructura principal - ángulo $\theta$
    \item \textbf{Gimbal ($g$):} Marco motorizado - ángulo $\alpha$
    \item \textbf{Flywheel ($f$):} Disco rotando a velocidad $\Omega$
\end{itemize}
\end{block}

\begin{center}
\fbox{\parbox{0.7\textwidth}{\centering \textit{[Diagrama de cuerpo libre]}}}
\end{center}

\begin{block}{Ecuaciones de Movimiento (linealizadas)}
\begin{equation*}
J_{eq}\ddot{\theta} - (M_{eq}L_{eq}g)\theta + H\dot{\alpha} = 0
\end{equation*}
\begin{equation*}
I_\alpha \ddot{\alpha} - H\dot{\theta} = \tau_\alpha
\end{equation*}
donde $H = I_{fy} \Omega$ es el momento angular del volante.
\end{block}
\end{frame}

%-------------------------------------------------------
\section{Diseño e Implementación}
%-------------------------------------------------------
\begin{frame}{Diseño Mecánico}{Prototipo CAD}
%-------------------------------------------------------

\begin{center}
\fbox{\parbox{0.8\textwidth}{\centering \textit{[Imagen: Diseño CAD completo]}}}
\end{center}

\vspace{0.3cm}

\begin{columns}
\column{0.5\textwidth}
\begin{block}{Sistema CMG}
\begin{itemize}
    \item Volante de inercia
    \item Motor DC alta velocidad
    \item Gimbal motorizado
\end{itemize}
\end{block}

\column{0.5\textwidth}
\fbox{\parbox{0.9\textwidth}{\centering \textit{[Imagen: Detalle CMG]}}}
\end{columns}
\end{frame}

%-------------------------------------------------------
\begin{frame}{Arquitectura del Sistema}{Integración Completa}
%-------------------------------------------------------

\begin{center}
\fbox{\parbox{0.85\textwidth}{\centering \textit{[Diagrama: Arquitectura completa]}}}
\end{center}

\vspace{0.3cm}

\begin{block}{Flujo de Control}
\begin{enumerate}
    \item IMU mide inclinación $\theta$ y velocidad angular
    \item Filtro complementario procesa señales
    \item Controlador PID calcula acción de control
    \item Actuador ajusta gimbal (ángulo $\alpha$)
    \item Torque giroscópico estabiliza la bicicleta
\end{enumerate}
\end{block}
\end{frame}

%-------------------------------------------------------
\begin{frame}{Sistema Electrónico}{Componentes}
%-------------------------------------------------------

\begin{columns}
\column{0.5\textwidth}
\begin{block}{Hardware}
\begin{itemize}
    \item ESP32 (controlador principal)
    \item MPU9250 (IMU 9-DOF)
    \item BTS7960 (driver motor DC)
    \item Batería LiFePO4
    \item BMS
\end{itemize}
\end{block}

\column{0.5\textwidth}
\begin{center}
\fbox{\parbox{0.9\textwidth}{\centering \textit{[Diagrama circuito]}}}
\end{center}
\end{columns}

\vspace{0.3cm}

\begin{center}
\fbox{\parbox{0.7\textwidth}{\centering \textit{[Imagen: Implementación electrónica]}}}
\end{center}
\end{frame}

%-------------------------------------------------------
\section{Resultados}
%-------------------------------------------------------
\begin{frame}{Prototipo Final}{Sistema Completo}
%-------------------------------------------------------

\begin{center}
\fbox{\parbox{0.8\textwidth}{\centering \textit{[Imagen: Prototipo ensamblado completo]}}}
\end{center}

\vspace{0.3cm}

\begin{columns}
\column{0.5\textwidth}
\fbox{\parbox{0.9\textwidth}{\centering \textit{[Imagen: Vista lateral]}}}

\column{0.5\textwidth}
\fbox{\parbox{0.9\textwidth}{\centering \textit{[Imagen: Detalle gimbal]}}}
\end{columns}
\end{frame}

%-------------------------------------------------------
\begin{frame}{Resultados Experimentales}{Desempeño del Sistema}
%-------------------------------------------------------

\begin{center}
\fbox{\parbox{0.85\textwidth}{\centering \textit{[Gráfica: Ángulo de inclinación vs Tiempo]}}}
\end{center}

\vspace{0.3cm}

\begin{columns}
\column{0.5\textwidth}
\begin{block}{Métricas}
\begin{itemize}
    \item Error en estado estacionario: $< 2°$
    \item Tiempo de asentamiento: [X] s
    \item Compensación de perturbaciones
\end{itemize}
\end{block}

\column{0.5\textwidth}
\fbox{\parbox{0.9\textwidth}{\centering \textit{[Gráfica: Señal de control]}}}
\end{columns}
\end{frame}

%-------------------------------------------------------
\begin{frame}{Validación}{Logros Alcanzados}
%-------------------------------------------------------

\begin{block}{Resultados Técnicos}
\begin{itemize}
    \item[$\checkmark$] Estabilización activa en reposo validada experimentalmente
    \item[$\checkmark$] Compensación efectiva de perturbaciones laterales
    \item[$\checkmark$] Control PID robusto y eficiente
    \item[$\checkmark$] Modelo matemático validado
\end{itemize}
\end{block}

\begin{center}
\fbox{\parbox{0.7\textwidth}{\centering \textit{[Imagen/Video: Demostración funcionamiento]}}}
\end{center}

\begin{alertblock}{Comparación con Estado del Arte}
Desempeño comparable a desarrollos internacionales, con plataforma de bajo costo implementada localmente.
\end{alertblock}
\end{frame}

%-------------------------------------------------------
\section{Conclusiones}
%-------------------------------------------------------
\begin{frame}{Conclusiones}{Técnicas y Académicas}
%-------------------------------------------------------

\begin{block}{Conclusiones Técnicas}
\begin{itemize}
    \item Se validó experimentalmente el principio de estabilización giroscópica mediante CMG
    \item El controlador PID demostró ser efectivo para el objetivo planteado
    \item El modelo matemático captura adecuadamente la dinámica del sistema
    \item Primera implementación física nacional exitosa
\end{itemize}
\end{block}

\begin{block}{Conclusiones Académicas}
\begin{itemize}
    \item Integración exitosa de teoría de control con implementación práctica
    \item Contribución al desarrollo tecnológico nacional
    \item Generación de conocimiento aplicado a movilidad sostenible
\end{itemize}
\end{block}
\end{frame}

%-------------------------------------------------------
\begin{frame}{Recomendaciones}{Trabajo Futuro}
%-------------------------------------------------------

\begin{block}{Mejoras al Sistema}
\begin{itemize}
    \item Implementar controladores avanzados (ADRC, LQR, SMC)
    \item Optimizar diseño mecánico para reducir peso total
    \item Mejorar eficiencia energética del sistema
\end{itemize}
\end{block}

\begin{block}{Extensiones del Proyecto}
\begin{itemize}
    \item Integrar sistema de control de dirección para navegación
    \item Evaluar desempeño en condiciones de movimiento
    \item Desarrollar sistema de navegación autónoma
\end{itemize}
\end{block}

\begin{block}{Aplicaciones Potenciales}
\begin{itemize}
    \item Sistema de asistencia para ciclistas urbanos
    \item Plataformas de movilidad personal inteligente
    \item Herramienta educativa en control de sistemas
\end{itemize}
\end{block}
\end{frame}

%-------------------------------------------------------
\begin{frame}{Impacto}{Social, Ambiental y Tecnológico}
%-------------------------------------------------------

\begin{columns}
\column{0.5\textwidth}
\begin{block}{Impacto Social}
\begin{itemize}
    \item Mayor seguridad para ciclistas
    \item Inclusión de usuarios con menor destreza
    \item Promoción de movilidad activa
\end{itemize}
\end{block}

\begin{block}{Impacto Ambiental}
\begin{itemize}
    \item Reducción de emisiones
    \item Transporte limpio
    \item Calidad del aire
\end{itemize}
\end{block}

\column{0.5\textwidth}
\begin{block}{Impacto Tecnológico}
\begin{itemize}
    \item Avance en control de sistemas dinámicos
    \item Desarrollo de prototipos funcionales
    \item Base para futuras investigaciones
\end{itemize}
\end{block}

\vspace{0.3cm}

\begin{center}
\fbox{\parbox{0.9\textwidth}{\centering \textit{[Imagen: Movilidad sostenible]}}}
\end{center}
\end{columns}
\end{frame}

%-------------------------------------------------------
% FINAL PAGE
%-------------------------------------------------------

\begin{frame}[plain]
  \begin{center}
  \Large \textbf{¡Gracias por su atención!}\\[1em]
  \normalsize
  \textbf{Contacto:}\\
  jhhernandezm@udistrital.edu.co\\
  fareyesr@udistrital.edu.co\\[1em]
  \textbf{Director:} Dr. Alexander Jiménez Triana\\
  Grupo de Investigación ORCA\\
  Universidad Distrital Francisco José de Caldas
  \end{center}
\end{frame}

\end{document}
